% Литобзор: сетевые методы анализа временных рядов (глава 1)
% Подключается в part1.tex вместо заглушки "Сетевые методы".
% Источники: population_recurrence_analysis.md (Varley & Sporns 2022 и др.)
% Конкретный метод и результаты — в секции \ref{sec:v2}.

\subsection{Сетевые методы анализа временных рядов}\label{sec:network_methods}

Методы снижения размерности и оценки размерности, рассмотренные выше, описывают популяционную активность как облако точек в многомерном пространстве, не используя информацию о порядке наблюдений во времени. Между тем темпоральная структура несёт самостоятельную информацию о динамических режимах системы. Отдельное направление анализа многомерных временных рядов составляют \textit{сетевые методы}, в которых временной ряд преобразуется в граф, а его свойства изучаются инструментами теории сетей~\cite{Varley2022}.

Существуют три основных способа построить граф из временного ряда. \textit{Графы рекуррентности} связывают моменты времени, в которые система возвращалась в близкие состояния фазового пространства; \textit{графы видимости} кодируют структуру локальных экстремумов сигнала; \textit{сети порядковых паттернов} отражают последовательность относительных порядков значений. Из них наиболее развит аппарат графов рекуррентности, восходящий к работе Эккманна, Кампхорста и Рюэля~\cite{Eckmann1987}, в которой граф рекуррентности был предложен как способ визуализации возвратов фазовой траектории. Введение количественных мер длин диагональных и вертикальных линий графа (recurrence quantification analysis) превратило метод в инструмент анализа детерминизма, ламинарности и нестационарности динамики~\cite{Marwan2007}.

Развитие теории сложных сетей дало дополнительное направление: матрицу рекуррентности можно трактовать как матрицу смежности комплексной сети и изучать её сетевыми мерами~--- так называемые сети рекуррентности~\cite{Donges2012}. Для совместного анализа нескольких сигналов разработаны многомерные расширения: кросс-рекуррентные и совместные (joint) графы рекуррентности~\cite{Romano2004}, а также мультиплексные сети рекуррентности~\cite{Eroglu2018}, в которых индивидуальные графы образуют слои многослойного сетевого объекта. Эти конструкции естественно ставят вопрос о применении графов рекуррентности к популяционным записям нейронной активности.

Несмотря на зрелость методологии, её применение к нейронным данным остаётся ограниченным. Подавляющее большинство работ использует графы рекуррентности для электрофизиологических сигналов (ЭЭГ, локальные потенциалы) или для отдельных временных рядов межспайковых интервалов. К сигналам кальциевой флуоресценции метод почти не применялся, а популяционные обобщения для записей отдельных нейронов до недавнего времени не предлагались. Развитие метода мультиплексного графа рекуррентности популяции и его применение к данным кальциевого имиджинга гиппокампа описаны в секции~\ref{sec:v2}.